\begin{table}[H]
\centering
\caption{Standardized Effect Sizes for Main Outcomes}
\label{tab:sde}
\begin{threeparttable}
\begin{adjustbox}{max width=\textwidth}
\begin{tabular}{llcccccc}
\toprule
Outcome & Specification & $\hat{\beta}$ & SE & SD($Y$) & SDE & SE(SDE) & Classification \\
\midrule
\multicolumn{8}{l}{\textit{Panel A: Pooled}} \\
Amihud illiquidity & Pooled DiD & -0.1074 & 0.0816 & 1.991 & -0.0539 & 0.0410 & Moderate negative \\
Turnover & Pooled DiD & 0.0977 & 0.0765 & 1.710 & 0.0571 & 0.0448 & Moderate positive \\
\midrule
\multicolumn{8}{l}{\textit{Panel B: Heterogeneous}} \\
Amihud (large treated) & Sample split & -0.1265 & 0.1000 & 1.991 & -0.0635 & 0.0502 & Moderate negative \\
Amihud (small treated) & Sample split & -0.0881 & 0.1032 & 1.991 & -0.0442 & 0.0518 & Small negative \\
\bottomrule
\end{tabular}
\end{adjustbox}
\end{threeparttable}
\par\vspace{0.3em}
{\footnotesize \textit{Notes:} \textbf{Country:} South Korea. \textbf{Research question:} Does mandatory English-language financial disclosure for large KOSPI firms reduce stock market illiquidity for treated firms relative to untreated KOSPI peers? \textbf{Policy mechanism:} Korea's FSC required firms with total assets at or above KRW 10 trillion to file financial statements in English starting January 2024, reducing language-based information barriers for foreign investors who previously could not read Korean-only filings. \textbf{Outcome definition:} Log of weekly Amihud illiquidity ratio (average absolute daily return divided by KRW trading volume in billions), where higher values indicate less liquid markets. \textbf{Treatment:} Binary: firms with total assets $\geq$ KRW 10 trillion mandated to file in English (Phase 1) vs.\ firms below the threshold. \textbf{Data:} Yahoo Finance daily OHLCV and FinanceDataReader firm characteristics for KOSPI-listed firms, January 2022--December 2025, at the firm-week level; 57,707 firm-week observations across 295 firms. \textbf{Method:} Two-way fixed effects DiD (firm and week FE), standard errors clustered at the firm level. \textbf{Sample:} Top 300 KOSPI firms by market capitalization; weeks with $\geq$3 trading days. Classification refers to magnitude, not statistical significance: Large ($|$SDE$| > 0.15$), Moderate ($0.05$--$0.15$), Small ($0.005$--$0.05$), Null ($< 0.005$).}
\end{table}
